\chapter{Trading dans un cloud public}\label{ch:nw:trading-in-a-public-cloud}

Dans la documentation d'un fournisseur de cloud, la \omterm{def:nw:trading-in-a-public-cloud:azid}{zone de disponibilité} appelée us-east-1a dans un compte n'est pas forcément la zone appelée us-east-1a
dans un autre~: «~AWS associe au hasard les \omterm{def:nw:trading-in-a-public-cloud:azid}{zones de disponibilité} physiques aux noms de \omterm{def:nw:trading-in-a-public-cloud:azid}{zones de disponibilité} de chaque compte AWS~». Seul
l'identifiant de zone, un code tel que use1-az1, désigne le même emplacement physique dans tous les comptes. Deux sociétés qui conviennent de se retrouver dans la «~zone a~» peuvent être
dans des bâtiments différents~; une société de trading qui place ses machines dans la zone que sa plateforme appelle «~a~» peut se trouver à une frontière de zone de celle-ci, et
le payer à chaque aller-retour.

La partie~IV de ce livre traite des plateformes qui fonctionnent dans un \omterm{def:nw:trading-in-a-public-cloud:cloud}{cloud public}, là où se fait l'essentiel du trading de crypto-actifs (chapitre~17), et du
cloud lui-même comme lieu depuis lequel trader. Ce chapitre décrit la géographie du cloud (régions, zones et identifiants qui les localisent),
les outils qui rapprochent les machines, le réseau et la virtualisation entre un programme et le câble, le bruit qu'ajoute
le matériel partagé, puis les horloges et les prix. Rien de tout cela n'est mesuré dans ce livre~: la documentation des fournisseurs est citée et datée, une
étude de mesure publiée fournit les allers-retours, et le reste relève d'une simulation signalée comme telle.

\section{Régions, zones de disponibilité et identifiants de zone}

\begin{definition}[Cloud public]\label{def:nw:trading-in-a-public-cloud:cloud}
Un \emph{cloud public}\index{cloud public} est un ensemble de calcul, de stockage et de réseau qu'un fournisseur exploite dans ses propres centres
de données et loue à tout client à la demande, à l'heure ou à l'unité, par une interface de programmation, sur un matériel partagé entre clients, sauf si
ceux-ci paient des machines dédiées.
\end{definition}

\begin{definition}[Zone de disponibilité, identifiant de zone de disponibilité]\label{def:nw:trading-in-a-public-cloud:azid}
Une \emph{zone de disponibilité}\index{zone de disponibilité} est l'un des groupes isolés de centres de données qui composent une région de cloud, dotés
de leur propre alimentation électrique, de leur propre refroidissement et de leur propre réseau, reliés aux autres zones de la région par des liaisons à faible latence. Son \emph{identifiant de zone de
disponibilité}\index{identifiant de zone de disponibilité} est le code que le fournisseur attribue à la zone physique, identique dans tous les comptes, à la différence du
nom de la zone, que le fournisseur peut attribuer différemment dans chaque compte.
\end{definition}

Une région de cloud (livre~3) est une zone géographique~; ses zones sont des bâtiments ou des campus situés à quelques kilomètres les uns des autres, choisis de sorte que
la panne d'une zone n'entraîne pas celle des autres. Pour une société de trading, la conséquence est celle du chapitre~9, sous une forme nouvelle~: deux machines d'une même zone
communiquent plus vite que deux machines de zones différentes, et le seul moyen de savoir dans quelle zone se trouvent les serveurs d'une plateforme est de passer par l'identifiant.

\begin{omfigure}
\begin{tikzpicture}[font=\footnotesize,b/.style={draw,rounded corners=2pt,minimum height=0.5cm,minimum width=1.6cm,align=center,inner sep=2pt},>=Stealth]
\node[font=\small] at (0,2.6) {compte A (une société)};
\node[font=\small] at (8.4,2.6) {compte B (une plateforme)};
\node[font=\small] at (4.2,2.6) {zones physiques};
\foreach \n/\y in {a/1.8,b/1.1,c/0.4,d/-0.3}{\node[b,fill=omDef!15] (A\n) at (0,\y) {us-east-1\n}; \node[b,fill=omDef!15] (B\n) at (8.4,\y) {us-east-1\n};}
\foreach \z/\y in {4/1.8,6/1.1,1/0.4,2/-0.3}{\node[b,fill=omThm!15] (Z\z) at (4.2,\y) {use1-az\z};}
\draw[->] (Aa) -- (Z4); \draw[->] (Ab) -- (Z6); \draw[->] (Ac) -- (Z1); \draw[->] (Ad) -- (Z2);
\draw[->] (Bb) -- (Z4); \draw[->] (Bc) -- (Z6); \draw[->] (Bd) -- (Z1); \draw[->] (Ba) -- (Z2);
\node[font=\scriptsize,align=center,omMeth] at (4.2,-1.1) {«~us-east-1a~» chez A est «~us-east-1b~» chez B~: les noms diffèrent, les identifiants non};
\end{tikzpicture}

\omcaption{Noms de zones et identifiants de zones pour deux comptes, d'après les données de test du chapitre (la forme de réponse documentée de
l'appel describe-availability-zones du fournisseur, avec des valeurs synthétiques~; quatre des six zones sont représentées). La société lit
l'identifiant de zone de la plateforme, et non son nom.}\label{fig:nw:trading-in-a-public-cloud:names}
\end{omfigure}

\begin{proposition}[Se fier aux noms]\label{prop:nw:trading-in-a-public-cloud:names}
Si une région compte $n$ zones et que les noms de chaque compte sont un mélange uniforme indépendant de ces zones, les zones de deux comptes qui
portent le même nom sont la même zone physique avec probabilité $1/n$. Si un aller-retour au sein d'une zone prend $t_s$ et un aller-retour entre zones $t_c$, choisir par nom plutôt que par
identifiant coûte $(1 - 1/n)(t_c - t_s)$ par aller-retour en moyenne.
\end{proposition}

\begin{proof}
Fixons la zone de la plateforme. La zone de la société portant le même nom est une zone physique tirée au hasard, elle est donc celle de la plateforme avec probabilité $1/n$~; le
reste du temps, l'aller-retour traverse des zones et prend $t_c$ au lieu de $t_s$.
\end{proof}

\omcode{code/firm/cloudplan/firm_cloudplan.py}{34}{56}{Noms de zones résolus en identifiants à partir d'une réponse describe-availability-zones, et
\cref{prop:nw:trading-in-a-public-cloud:names}.}\label{lst:nw:trading-in-a-public-cloud:zones}

Avec les allers-retours publiés de la quatrième section de ce chapitre (médiane de \qty{270}{\micro\second} dans une zone, \qty{555}{\micro\second} entre zones),
se fier aux noms coûte \qty{190}{\micro\second} par aller-retour en moyenne dans une région de trois zones, et \qty{238}{\micro\second} dans une région de six
(\cref{fig:nw:trading-in-a-public-cloud:penalty})~: bien plus que n'importe quel réglage de la machine ne permettrait de récupérer.

\begin{omfigure}
\begin{tikzpicture}
\begin{axis}[ybar,width=10cm,height=5.4cm,bar width=16pt,ymin=0,ymax=280,xtick=data,xlabel={zones de la région},
  ylabel={pénalité espérée (\si{\micro\second})},tick label style={font=\footnotesize},label style={font=\small},grid=major,
  grid style={gray!25},nodes near coords,nodes near coords style={font=\scriptsize},table/col sep=comma,enlarge x limits=0.2]
\addplot[fill=omDef!55,draw=omDef] table[x=n,y=penalty_us]{figdata/networks/16-trading-in-a-public-cloud/penalty.csv};
\end{axis}
\end{tikzpicture}

\omcaption{Pénalité espérée par aller-retour lorsqu'on choisit une zone par nom plutôt que par identifiant (\cref{prop:nw:trading-in-a-public-cloud:names}),
en fonction du nombre de zones de la région, avec les allers-retours médians publiés au sein d'une zone et entre zones. Données~: \texttt{fig\_cloud.py},
\texttt{nw\_cloud.penalty\_rows()}.}\label{fig:nw:trading-in-a-public-cloud:penalty}
\end{omfigure}

La même logique régionale touche les plateformes traditionnelles~: CME et Google Cloud construisent une région de cloud privée et un centre de
\omterm{def:nw:colocation-products-and-how-they-are-sold:colo}{colocation} à Aurora (chapitre~10), si bien que les zones et les identifiants pourraient un jour compter pour les contrats à terme autant qu'ils comptent déjà pour les crypto-actifs.

\section{Groupes de placement et choix des instances}

\begin{definition}[Groupe de placement]\label{def:nw:trading-in-a-public-cloud:pg}
Un \emph{groupe de placement}\index{groupe de placement} est une instruction du client au fournisseur de cloud sur l'emplacement d'un ensemble de
ses machines les unes par rapport aux autres~: rapprochées pour une faible latence (un cluster), éloignées pour la résilience (un spread), ou placées sur un matériel aux capacités particulières.
\end{definition}

Les fournisseurs nomment l'outil différemment et promettent des choses différentes. Les \omterm{def:nw:trading-in-a-public-cloud:pg}{groupes de placement} de type cluster d'AWS maintiennent les instances dans une même \omterm{def:nw:trading-in-a-public-cloud:azid}{zone de disponibilité},
dans un même segment du réseau à forte bande passante de bissection, et les recommande pour une faible latence~; ses \omterm{def:nw:trading-in-a-public-cloud:pg}{groupes de placement} de temps de précision placent
les instances sur un matériel ayant un accès direct à des sources de temps de haute précision. Les proximity placement groups d'Azure rendent les ressources «~physiquement situées
à proximité les unes des autres~». Les compact placement policies de Google Cloud placent les instances à proximité les unes des autres dans une zone «~selon le principe du meilleur effort~», éventuellement
dans une distance maximale. Aucun ne publie de chiffre de latence, ce qui explique que la société mesure (chapitre~17).

\begin{dated}{2026-09}{Placement et services de temps, d'après la documentation des fournisseurs}\label{dat:nw:trading-in-a-public-cloud:docs}
\begin{itemize}
\item AWS~: noms de zones associés au hasard par compte, identifiants de zones cohérents d'un compte à l'autre~; stratégies de placement cluster, partition, spread et
  temps de précision~; les groupes de temps de précision donnent une source NTP locale renforcée, une horloge matérielle PTP pour une synchronisation précise à la microseconde (Linux),
  et un horodatage matériel des paquets avec une résolution à la nanoseconde, sans frais supplémentaires.
\item Réseau AWS~: réseau amélioré par SR-IOV~; ENA Express, qui porte la bande passante d'un flux unique de 5 à 25~Gb/s et réduit la
  latence de queue au sein d'une zone~; \omterm{def:nw:trading-in-a-public-cloud:bare}{instances bare-metal} sans surcharge de virtualisation.
\item Azure~: proximity placement groups. Google Cloud~: compact placement policies, au meilleur effort, avec une distance maximale facultative.
\end{itemize}
\end{dated}

Le choix des instances suit la même logique que celui des serveurs du chapitre~8~: l'horloge et le nombre de cœurs de la machine, sa bande passante réseau, et le fait que
la \omterm{def:nw:networking-for-trading:frame}{carte réseau} et ses déchargements atteignent le programme sans qu'un hyperviseur fasse obstacle. Dans le cloud, le choix se fait sur une ligne d'un catalogue assorti
d'un prix horaire, et la société peut en essayer plusieurs et garder le meilleur (la loterie des instances du chapitre~17).

\section{Cartes réseau, virtualisation et bare metal}

\begin{definition}[Virtualisation d'E/S à racine unique]\label{def:nw:trading-in-a-public-cloud:sriov}
La \emph{virtualisation d'E/S à racine unique}\index{virtualisation d'E/S à racine unique} (single-root I/O virtualisation, SR-IOV) permet à une \omterm{def:nw:networking-for-trading:frame}{carte réseau} physique de présenter plusieurs fonctions virtuelles,
chacune étant assignée directement à une machine virtuelle, de sorte que le pilote de la machine dialogue avec le matériel sans que l'hyperviseur recopie ses
paquets.
\end{definition}

\begin{definition}[Instance bare-metal]\label{def:nw:trading-in-a-public-cloud:bare}
Une \emph{instance bare-metal}\index{instance bare-metal} (bare-metal instance) est un serveur physique entier loué par l'intermédiaire de l'interface du cloud, sans hyperviseur
entre le système d'exploitation du client et le matériel, et sans autre client sur la machine.
\end{definition}

Entre un programme de trading sur une instance cloud et le câble, on trouve, dans l'ordre~: le programme, le noyau ou une bibliothèque de contournement du noyau (kernel bypass) (livre~13),
la fonction virtuelle d'une carte SR-IOV, la \omterm{def:nw:networking-for-trading:frame}{carte réseau} propre au fournisseur et son logiciel, le réseau physique de l'hôte, et la
structure du réseau de la région. Le fournisseur ne documente le milieu de cette pile qu'à grands traits~: mise en réseau améliorée par SR-IOV, ENA Express sur un
transport propriétaire pour un débit plus élevé par flux unique et une latence de queue plus faible au sein d'une zone, et le bare metal pour les clients qui ont besoin du matériel pour eux seuls.
La société peut contrôler les extrémités~: son programme et ses réglages du noyau (chapitre~3), ainsi que l'instance, le placement et la zone qu'elle loue.

\section{Gigue et voisins bruyants}

\begin{definition}[Temps volé au CPU, voisin bruyant]\label{def:nw:trading-in-a-public-cloud:steal}
Le \emph{temps volé au CPU}\index{temps volé au CPU} (CPU steal time) est le temps pendant lequel un CPU virtuel était prêt à s'exécuter, mais où l'hyperviseur a fait tourner autre chose sur le CPU physique~; Linux
le signale comme «~attente involontaire~». Un \emph{voisin bruyant}\index{voisin bruyant} (noisy neighbour) est un autre locataire du même serveur physique ou du même réseau, dont la charge
prélève des ressources (CPU, cache, bande passante mémoire, files d'attente réseau) sur la machine d'un client et lui ajoute ainsi du délai.
\end{definition}

La seule mesure publiée sur laquelle s'appuie ce chapitre, celle de Hilyard et de ses coauteurs, a enregistré les allers-retours entre de petites machines virtuelles dans une
région US East, chez trois fournisseurs, pendant six heures. Sur AWS, les machines d'un même sous-réseau, dans une même zone, présentaient un aller-retour médian de \qty{270}{\micro\second}, un
99\up{e}~centile de \qty{395}{\micro\second} et un 99.99\up{e}~centile de \qty{760}{\micro\second}~; les machines situées dans des zones différentes affichaient 555, 755 et
\qty{1905}{\micro\second}~; et les pires allers-retours atteignaient 2\,900 fois la moyenne. Le modèle ajuste une loi lognormale à chaque médiane et à chaque 99e
centile et ajoute des événements de queue rares calibrés sur le 99.99\up{e}~centile (\cref{tab:nw:trading-in-a-public-cloud:rtt})~; un troisième cas, un hôte partagé avec un
\omterm{def:nw:trading-in-a-public-cloud:steal}{voisin bruyant}, conserve le corps de la distribution de la même zone et ajoute des pics de 5 à 50 fois un aller-retour sur 500, ce qui relève d'une hypothèse.

\begin{table}[ht]
\centering\small
\begin{tabular}{lrrrrr}
\toprule
Cas & p50 & p90 & p99 & p99.9 & p99.99 \\
\midrule
Même zone (publié) & 270 & -- & 395 & -- & 760 \\
Même zone (modèle) & 270 & 333 & 396 & 463 & 743 \\
Zones différentes (publié) & 555 & -- & 755 & -- & 1\,905 \\
Zones différentes (modèle) & 555 & 657 & 755 & 851 & 1\,908 \\
Hôte partagé, \omterm{def:nw:trading-in-a-public-cloud:steal}{voisin bruyant} (modèle) & 270 & 333 & 400 & 7\,474 & 13\,940 \\
\bottomrule
\end{tabular}
\caption{Centiles des allers-retours en microsecondes~: la mesure publiée (Hilyard et al., AWS, même sous-réseau et zones différentes) et la simulation
que \texttt{firm.cloudplan} ajuste à celle-ci~; le cas du \omterm{def:nw:trading-in-a-public-cloud:steal}{voisin bruyant} est une hypothèse. Données~: \texttt{nw\_cloud.percentiles()}.}\label{tab:nw:trading-in-a-public-cloud:rtt}
\end{table}

\begin{omfigure}
\begin{tikzpicture}
\begin{semilogyaxis}[width=11cm,height=6cm,xmin=-0.2,xmax=4.2,ymin=150,ymax=30000,xtick={0,1,2,3,4},
  xticklabels={p50,p90,p99,p99.9,p99.99},ylabel={aller-retour (\si{\micro\second}, log)},tick label style={font=\footnotesize},
  label style={font=\small},grid=major,grid style={gray!25},legend pos=north west,legend style={font=\footnotesize},table/col sep=comma]
\addplot[omDef,thick,mark=*,mark size=1.5pt] table[x=pos,y=same]{figdata/networks/16-trading-in-a-public-cloud/rtt.csv};
\addplot[omThm,thick,mark=square*,mark size=1.5pt] table[x=pos,y=cross]{figdata/networks/16-trading-in-a-public-cloud/rtt.csv};
\addplot[omMeth,thick,dashed,mark=triangle*,mark size=1.8pt] table[x=pos,y=noisy]{figdata/networks/16-trading-in-a-public-cloud/rtt.csv};
\legend{même zone,zones différentes,{hôte partagé, voisin bruyant}}
\end{semilogyaxis}
\end{tikzpicture}

\omcaption{Simulation ajustée à une mesure publiée~: centiles de l'aller-retour au sein d'une zone et entre zones, et sur un hôte partagé avec un
\omterm{def:nw:trading-in-a-public-cloud:steal}{voisin bruyant} supposé. Franchir une zone double le corps de la distribution~; un \omterm{def:nw:trading-in-a-public-cloud:steal}{voisin bruyant} laisse le corps intact et multiplie la queue. Données~:
\texttt{fig\_cloud.py}.}\label{fig:nw:trading-in-a-public-cloud:rtt}
\end{omfigure}

\omcode{code/firm/cloudplan/firm_cloudplan.py}{68}{74}{Le modèle d'aller-retour~: une loi lognormale ajustée à une médiane et à un 99\up{e}~centile, avec des pics rares.}\label{lst:nw:trading-in-a-public-cloud:rtt}

La leçon est celle du livre~13, dans un cadre plus rude~: la médiane n'est pas le problème. Une stratégie qui court dans le cloud se joue dans la
queue, et la queue appartient aux autres locataires, sauf si la société loue la machine entière~; le 99.9\up{e}~centile du cas du \omterm{def:nw:trading-in-a-public-cloud:steal}{voisin bruyant} vaut 16 fois
la valeur de la même zone. Le bare metal supprime les voisins sur l'hôte, pas ceux du réseau.

\section{Horloges et coûts dans le cloud}

Les horloges du chapitre~4 ont une version pour le cloud. AWS propose un service de temps local disponible sur chaque instance et, dans les \omterm{def:nw:trading-in-a-public-cloud:pg}{groupes de placement} à temps de précision,
une horloge matérielle PTP pour une synchronisation «~précise à la microseconde~» et des horodatages matériels des paquets avec une résolution nanoseconde~; cela suffit à
horodater les allers-retours du chapitre lui-même, ce qu'une société doit faire pour voir sa queue.

\begin{definition}[Frais de transfert de données]\label{def:nw:trading-in-a-public-cloud:dtc}
Les \emph{frais de transfert de données}\index{frais de transfert de données} (data-transfer charge) sont la commission que facture le fournisseur de cloud pour les données qui franchissent une frontière de son réseau (entre zones,
entre régions ou vers l'internet), par gigaoctet et par sens, en plus du prix des machines.
\end{definition}

\begin{dated}{2026-09}{Prix publiés du cloud}\label{dat:nw:trading-in-a-public-cloud:prices}
Prix à la demande AWS pour Linux dans US East (N. Virginia), d'après la grille tarifaire publiée le 25 septembre 2026~: c7i.xlarge (4 vCPU) 0.1785~USD de l'heure~;
c7i.4xlarge (16 vCPU, jusqu'à 12.5 Gb/s) 0.714~; c6in.4xlarge (16 vCPU, jusqu'à 50 Gb/s) 0.9072~; c7i.24xlarge et le bare-metal c7i.metal-24xl (96
vCPU, 37.5 Gb/s) 4.284 chacun~; c7i.metal-48xl 8.568. Les données au sein d'une même zone sont gratuites~; les données qui franchissent les zones sont facturées 0.01~USD par gigaoctet dans chaque
sens.
\end{dated}

\begin{omfigure}
\begin{tikzpicture}
\begin{axis}[xbar stacked,width=11cm,height=4.8cm,bar width=10pt,xmin=0,xmax=7200,ytick=data,
  yticklabels={quatre c7i.4xlarge dans une zone,quatre c7i.4xlarge sur deux zones,deux c7i.metal-24xl dans une zone},y dir=reverse,
  xlabel={coût mensuel (USD)},tick label style={font=\footnotesize,/pgf/number format/1000 sep={\,}},label style={font=\small},grid=major,grid style={gray!25},
  legend style={font=\footnotesize,at={(0.5,-0.4)},anchor=north,draw=none},legend columns=2,table/col sep=comma]
\addplot[fill=omDef!55,draw=omDef] table[y expr=\coordindex,x=compute]{figdata/networks/16-trading-in-a-public-cloud/plans.csv};
\addplot[fill=omThm!45,draw=omThm] table[y expr=\coordindex,x=transfer]{figdata/networks/16-trading-in-a-public-cloud/plans.csv};
\legend{instances (730 heures),données entre zones (20 To dans chaque sens)}
\end{axis}
\end{tikzpicture}

\omcaption{Trois formules mensuelles chiffrées d'après le catalogue daté de l'\cref{dat:nw:trading-in-a-public-cloud:prices}~: quatre instances de 16 vCPU dans une même zone,
la même répartition sur deux zones échangeant 20 téraoctets par mois, et deux serveurs bare-metal dans une même zone. Données~: \texttt{fig\_cloud.py},
\texttt{nw\_cloud.plan\_costs()}.}\label{fig:nw:trading-in-a-public-cloud:plans}
\end{omfigure}

Les prix sont modestes à côté de la \omterm{def:nw:colocation-products-and-how-they-are-sold:colo}{colocation} du chapitre~9~: quatre instances à 16 cœurs coûtent 2\,085~USD par mois, deux serveurs complets à 96 cœurs 6\,255~USD, là où une seule
baie Nasdaq coûte 7\,230~USD avant l'alimentation et les connexions. Répartir les quatre instances sur deux zones et échanger 20 téraoctets par mois ajoute
400~USD de \omterm{def:nw:trading-in-a-public-cloud:dtc}{frais de transfert de données}, et \qty{285}{\micro\second} à chaque médiane d'aller-retour entre elles. Dans le cloud, ce qui coûte cher n'est pas la
machine mais la distance.

\begin{method}[Placement d'un système de trading dans une région cloud]\label{met:nw:trading-in-a-public-cloud:place}
\begin{enumerate}
\item Trouver la région de la plateforme et, si elle la publie, l'identifiant de zone de ses serveurs (chapitres~17 et~18)~; sinon, le trouver par mesure.
\item Résoudre les noms de zone en identifiants dans chaque compte utilisé, et placer par identifiant.
\item Regrouper les machines sensibles à la latence dans une même zone, dans un \omterm{def:nw:trading-in-a-public-cloud:pg}{groupe de placement} de type cluster (et, si l'on a besoin de l'horloge, de type précision temporelle)~; envisager le
  bare-metal pour celles qui se disputent la course.
\item Horodater ses propres allers-retours avec le service de temps du fournisseur et mesurer la queue de distribution, non la médiane~; refaire la mesure après chaque maintenance
  ou toute relance.
\item Chiffrer les machines et le transfert de données ensemble~; réserver le trafic inter-zones à ce qui exige de la résilience.
\end{enumerate}
\end{method}

\section{Tutoriel~: zones, queues de distribution et facture}\label{tut:nw:trading-in-a-public-cloud:tut}

\begin{tutorial}
\textbf{Objectif.} Résoudre les noms de zone, simuler des allers-retours à partir d'une mesure publiée et chiffrer trois formules~; aucun compte cloud n'est nécessaire.
\textbf{État final~:} \cref{fig:nw:trading-in-a-public-cloud:penalty,fig:nw:trading-in-a-public-cloud:rtt,fig:nw:trading-in-a-public-cloud:plans}
et le \cref{tab:nw:trading-in-a-public-cloud:rtt}.
\begin{tutsteps}
\item \textbf{Zones.} \texttt{firm\_cloudplan.load\_zones} lit les deux jeux de données dans \texttt{data/} (la forme de réponse documentée, valeurs
  synthétiques)~; \texttt{same\_physical} et \texttt{locate} les comparent (voir le \cref{lst:nw:trading-in-a-public-cloud:zones}).
\item \textbf{Pénalité.} \texttt{nw\_cloud.penalty\_rows()} évalue la \cref{prop:nw:trading-in-a-public-cloud:names} pour 2 à 6 zones.
\item \textbf{Allers-retours.} \texttt{sample\_rtt} (voir le \cref{lst:nw:trading-in-a-public-cloud:rtt}) et les trois \texttt{CASES}.
\item \textbf{Facture.} \texttt{load\_catalogue} et \texttt{monthly\_cost} chiffrent \texttt{PLANS}.
\end{tutsteps}
\textbf{Pour aller plus loin.} Exécuter le code du chapitre sur de vraies instances dans deux zones, avec des horodatages matériels, remplacer le cas ajusté
par vos propres percentiles, et comparer.
\end{tutorial}

\section{Construction~: le plan de placement cloud}\label{bld:nw:trading-in-a-public-cloud:cloudplan}

\begin{build}
\bfield{Rôle} Où se trouvent les machines de la société dans une région cloud, ce qu'elles coûtent et à quoi ressemblent leurs allers-retours~: la base des chapitres~17 à~19
et des lignes cloud du plan du chapitre~29.
\bfield{Interface} \texttt{firm\_cloudplan}~: \texttt{load\_zones}, \texttt{same\_physical}, \texttt{locate}, \texttt{p\_same\_name\_same\_zone},
\texttt{expected\_penalty\_us}, \texttt{RTT\_PUBLISHED}, \texttt{Rtt}, \texttt{sample\_rtt}, \texttt{Instance}, \texttt{load\_catalogue},
\texttt{monthly\_cost}.
\bfield{Règles} Placer par identifiant de zone, jamais par nom~; les prix viennent du catalogue daté, avec leur source~; les allers-retours simulés indiquent
à quelle mesure ils sont ajustés.
\bfield{Tests d'acceptation} \texttt{code/firm/cloudplan/tests/}~: les noms et identifiants des jeux de données, la probabilité $1/n$ comparée à une
simulation de permutation, la médiane et le 99\up{e}~percentile de l'ajustement, et une facture calculée à la main.
\bfield{Pour aller plus loin} Lire les réponses en direct et les grilles tarifaires par les interfaces du fournisseur~; ajouter des régions et des allers-retours inter-régions.
\end{build}

\begin{omsources}
\item Documentation AWS~: Availability Zone IDs~; placement strategies et precision time placement groups~; enhanced networking, ENA Express et
Nitro bare-metal instances~; Amazon VPC pricing et la liste de prix publique (septembre 2026).
\item Microsoft Learn, proximity placement groups~; Google Cloud, placement policies.
\item O.~Hilyard, B.~Cui, M.~Webster, A.~Bangalore Muralikrishna et A.~Charapko, «~Cloudy Forecast: How Predictable is Communication Latency in the
Cloud?~», arXiv:2309.13169 (2023).
\end{omsources}

\section{Exercices}

\begin{exercise}[$\star$]\label{exo:nw:trading-in-a-public-cloud:1}
Dans une région de quatre zones, quelle est la probabilité que votre «~zone b~» soit la «~zone b~» de la plateforme~?
\end{exercise}

\begin{exercise}[$\star$]\label{exo:nw:trading-in-a-public-cloud:2}
Combien coûte par mois le déplacement de 5 téraoctets dans chaque sens entre deux zones~?
\end{exercise}

\begin{exercise}[$\star$]\label{exo:nw:trading-in-a-public-cloud:3}
Que mesure le \omterm{def:nw:trading-in-a-public-cloud:steal}{temps volé au CPU} (steal time), et pourquoi un programme de trading s'en soucie-t-il~?
\end{exercise}

\begin{exercise}[$\star\star$]\label{exo:nw:trading-in-a-public-cloud:4}
Pourquoi le passage d'une zone à une autre double-t-il la médiane de l'aller-retour dans la mesure publiée, alors qu'un \omterm{def:nw:trading-in-a-public-cloud:steal}{voisin bruyant} la laisse inchangée~?
\end{exercise}

\begin{exercise}[$\star\star$]\label{exo:nw:trading-in-a-public-cloud:5}
Que supprime une \omterm{def:nw:trading-in-a-public-cloud:bare}{instance bare-metal}, et que ne supprime-t-elle pas~?
\end{exercise}

\begin{exercise}[$\star\star$]\label{exo:nw:trading-in-a-public-cloud:6}
Comment trouver l'identifiant de zone des serveurs d'une plateforme qui ne le publie pas~?
\end{exercise}

\begin{exercise}[$\star\star\star$]\label{exo:nw:trading-in-a-public-cloud:7}
\emph{Programmation.} Avec \texttt{firm.cloudplan}, trouver la probabilité de pic à partir de laquelle le 99.9\up{e}~percentile du cas du \omterm{def:nw:trading-in-a-public-cloud:steal}{voisin bruyant} dépasse pour la première fois celui du cas
inter-zones.
\end{exercise}

\begin{exercise}[$\star\star\star$]\label{exo:nw:trading-in-a-public-cloud:8}
\emph{Trouver la faille.} «~Nous avons lancé nos machines dans us-east-1a, comme l'indique la plateforme dans sa documentation, donc nous sommes dans sa zone.~»
\end{exercise}

\section{Problème~: us-east-1a n'est pas votre us-east-1a}

\begin{problem}[{Problème du week-end --- le prix de la confiance accordée à un nom}]\label{pb:nw:trading-in-a-public-cloud:1}
Une plateforme fait tourner son moteur d'appariement (matching engine) dans une zone d'une région à six zones et communique à ses utilisateurs le nom de la zone dans son propre compte. Une société lance ses machines
dans la zone de même nom de son propre compte. Utiliser les allers-retours publiés~: \qty{270}{\micro\second} au sein d'une zone, \qty{555}{\micro\second} entre zones.

\medskip\noindent\textbf{Partie I --- Les noms.}
\begin{enumerate}
\item Avec des permutations uniformes indépendantes, quelle est la probabilité que la société se trouve dans la zone de la plateforme~?
\item Quelle est la pénalité moyenne d'aller-retour, et qu'en serait-il si la région comptait trois zones~?
\item Dans les jeux de données du chapitre, lesquels des noms du compte A correspondent physiquement à ceux du compte B~?
\item Sous quel nom le compte B voit-il le us-east-1a du compte A~?
\end{enumerate}

\medskip\noindent\textbf{Partie II --- La queue de distribution.}
\begin{enumerate}[resume]
\item Quels sont les percentiles 99\up{e} et 99.99\up{e} modélisés, au sein d'une zone et entre zones~?
\item À quel point le modèle est-il proche des percentiles 99.99\up{e} publiés~?
\item Que fait un \omterm{def:nw:trading-in-a-public-cloud:steal}{voisin bruyant} au 99.9\up{e}~percentile~?
\item Qu'est-ce qui compte le plus pour une stratégie de course, la zone ou le voisin~?
\end{enumerate}

\medskip\noindent\textbf{Partie III --- La facture.}
\begin{enumerate}[resume]
\item Combien coûtent quatre instances c7i.4xlarge par mois~?
\item Qu'ajoutent 20 téraoctets par mois dans chaque sens entre zones~?
\item Combien coûtent deux \omterm{def:nw:trading-in-a-public-cloud:bare}{instances bare-metal} c7i.metal-24xl, et à quel chiffre de \omterm{def:nw:colocation-products-and-how-they-are-sold:colo}{colocation} faut-il les comparer~?
\item Qu'est-ce qui coûte le moins cher, une erreur de zone ou une mise à niveau bare-metal, et pourquoi est-ce la mauvaise question~?
\end{enumerate}

\medskip\noindent\textbf{Partie IV --- Le verdict.}
\begin{enumerate}[resume]
\item Énoncer le \emph{résultat nommé}~: la probabilité que la zone de même nom de deux comptes soit la même zone physique lorsque les noms sont mélangés sur $n$
  zones, et la pénalité moyenne d'aller-retour liée au fait de faire confiance aux noms plutôt qu'aux identifiants.
\item Comment la société pourrait-elle vérifier son placement sans l'aide de la plateforme~?
\item Pourquoi les fournisseurs mélangent-ils les noms~?
\item Que devrait publier une plateforme, et que devrait demander une société~?
\item Quelles parties de ce problème relèvent de la documentation, lesquelles de la mesure, lesquelles de l'hypothèse~?
\item À quelle fréquence faut-il revérifier le placement~?
\item Le même problème existe-t-il dans une salle de \omterm{def:nw:colocation-products-and-how-they-are-sold:colo}{colocation}~?
\item En une phrase~: qu'apprend un nom de zone~?
\end{enumerate}
\end{problem}

\section{Questions d'entretien}

\begin{interviewq}[$\star$ \iqroles{développeur}]\label{iq:nw:trading-in-a-public-cloud:1}
Qu'est-ce qu'une \omterm{def:nw:trading-in-a-public-cloud:azid}{zone de disponibilité}, et pourquoi compte-t-elle pour la latence~?
\end{interviewq}

\begin{interviewq}[$\star\star$ \iqroles{développeur}]\label{iq:nw:trading-in-a-public-cloud:2}
La latence médiane de votre système de trading sur le cloud est correcte, mais son 99.9\up{e}~centile est dix fois pire. Où regardez-vous~?
\end{interviewq}

\begin{interviewq}[$\star\star$ \iqroles{développeur}]\label{iq:nw:trading-in-a-public-cloud:3}
Comment placez-vous les machines le plus près possible d'une plateforme qui tourne dans la même région cloud~?
\end{interviewq}

\begin{interviewq}[$\star\star$ \iqroles{développeur, chercheur}]\label{iq:nw:trading-in-a-public-cloud:4}
Pouvez-vous horodater les paquets avec précision dans un \omterm{def:nw:trading-in-a-public-cloud:cloud}{cloud public}~? Comment~?
\end{interviewq}

\begin{interviewq}[$\star\star$ \iqroles{développeur}]\label{iq:nw:trading-in-a-public-cloud:5}
Combien coûte l'exploitation d'un système de trading dans le cloud comparée à la \omterm{def:nw:colocation-products-and-how-they-are-sold:colo}{colocation}, et d'où viennent les coûts cachés~?
\end{interviewq}

\begin{interviewq}[$\star\star\star$ \iqroles{développeur, trader}]\label{iq:nw:trading-in-a-public-cloud:6}
Feriez-vous tourner une stratégie actions sensible à la latence depuis un \omterm{def:nw:trading-in-a-public-cloud:cloud}{cloud public}~? Que faudrait-il qu'il soit vrai~?
\end{interviewq}
